\chapter{ESTADO DEL ARTE}
\label{ch:estado_arte}

Este capítulo presenta una revisión general de las principales tecnologías
utilizadas en robots móviles autónomos para interiores. El análisis se centra
en las alternativas disponibles para plataformas de software, percepción,
mapeo, navegación, exploración y simulación, considerando su capacidad de
integración entre sí, sus requerimientos computacionales y su nivel de
madurez.

\section{Panorama general de la robótica móvil}
\label{sec:panorama_robotica_movil}

El desarrollo actual de robótica móvil abarca aplicaciones en dominios tan diversos como la logística, la agricultura y las flotas multirrobot, entre otros \citep{macenski2023desks}. La diversidad de aplicaciones ha dado lugar a
robots con distintas arquitecturas de locomoción. Las plataformas diferenciales son habituales en robots móviles pequeños, por su simplicidad mecánica y su capacidad de girar sobre sí mismas \parencites[Sec.~4.1]{corke2023}{ferreira2023differential}. Los
vehículos con dirección tipo Ackermann se emplean cuando se requiere un
movimiento semejante al de un automóvil. Las bases omnidireccionales
permiten desplazamientos laterales en espacios reducidos. Cada configuración ofrece
distinta maniobrabilidad y complejidad mecánica. Por ejemplo, un vehículo tipo
automóvil no puede girar sobre sí mismo \citep[Sec.~4.1]{corke2023}.

También existen diferencias en el nivel de autonomía. Algunos sistemas siguen
rutas previamente definidas, mientras que otros estiman su posición,
construyen mapas, planifican trayectorias y reaccionan ante cambios del
entorno. Por ejemplo, un robot logístico puede desplazarse entre estaciones
conocidas; un robot de inspección puede recorrer una instalación y registrar
mediciones; y una plataforma de investigación puede explorar un recinto
desconocido mientras genera su mapa.
Desde el punto de vista computacional, estas funciones pueden concentrarse en
un solo computador o distribuirse entre varios dispositivos. El control de
los actuadores suele requerir una ejecución rápida y periódica, mientras que
la percepción, el mapeo y la planificación demandan mayor capacidad de
procesamiento. Por ello, los sistemas robóticos suelen descomponerse en componentes
independientes que pueden ejecutarse en un mismo proceso, en varios procesos o en
varios computadores \citep{macenski2022ros2}.

\section{Plataformas de software y middleware robótico}
\label{sec:alternativas_middleware}

La integración de sensores, actuadores y algoritmos puede realizarse mediante
una plataforma robótica de propósito general, un middleware de comunicación o
una arquitectura desarrollada específicamente para el sistema. Entre las
alternativas más conocidas se encuentran ROS~1, ROS~2, YARP, Orocos RTT y
LCM, además de soluciones propietarias utilizadas en plataformas comerciales.
Estas alternativas no ofrecen exactamente el mismo nivel de funcionalidad.
Algunas proporcionan bibliotecas, herramientas de visualización, gestión de
paquetes y controladores de hardware, mientras que otras se concentran en la
comunicación entre procesos o en la ejecución de componentes de tiempo real.
La Tabla~\ref{tab:comparativa_middleware} presenta una comparación general.

\begin{table}[htbp]
\centering
\caption{Comparación general de plataformas y middleware para robótica.}
\label{tab:comparativa_middleware}
\small
\begin{tabular}{P{2.1cm} P{5.4cm} P{6.3cm}}
\toprule
\textbf{Alternativa} &
\textbf{Características generales} &
\textbf{Alcance declarado por el proyecto} \\
\midrule

ROS~1 &
Amplio conjunto histórico de paquetes, controladores y herramientas para
desarrollo robótico. &
Cuenta con una comunidad y base de software extensa, pero su última
distribución oficial finalizó su período de soporte en 2025.
\citep{ros_noetic_eol} \\

ROS~2 &
Plataforma modular para sistemas distribuidos, con diferentes mecanismos de
comunicación, calidad de servicio y herramientas para integrar hardware y
algoritmos. &
Ofrece un ecosistema amplio, aunque su capa de comunicaciones puede requerir
ajustes, por ejemplo para transmitir mensajes de gran tamaño.
\citep{macenski2022ros2} \\

YARP &
Middleware orientado a la comunicación entre módulos y a la abstracción de
dispositivos; se utiliza en plataformas que requieren conectar componentes
distribuidos. &
El proyecto se define en torno a la comunicación entre módulos y la
abstracción de dispositivos, sin declarar una pila propia de navegación
móvil \citep{yarp_docs}. \\

Orocos RTT &
Entorno basado en componentes para control robótico y aplicaciones con
requerimientos de tiempo real. &
El proyecto se orienta al control determinista por componentes y no
distribuye módulos de mapeo ni de planificación de trayectorias
\citep{orocos_rtt}.\\

LCM &
Sistema ligero de publicación y suscripción, orientado al intercambio de
mensajes con baja latencia. &
El proyecto se define como una biblioteca de mensajería de baja latencia y
no incluye componentes de percepción, mapeo ni navegación
\citep{lcm_docs}.\\

\bottomrule
\end{tabular}
\end{table}
En consecuencia, la selección de una plataforma de software depende de la
disponibilidad de controladores, los requerimientos temporales, la capacidad
computacional, la facilidad de integración, la comunidad disponible y el
grado de desarrollo propio que puede asumir el proyecto. La alternativa
adoptada para el prototipo y su relación con los restantes componentes se
presentan en el Capítulo~\ref{ch:solucion}.

\section{Percepción y mapeo del entorno}
\label{sec:sensores}

En robots de interiores, las principales alternativas de percepción incluyen
sensores \gls{LIDAR} 2D, cámaras RGB-D, sistemas de visión estéreo y sensores
\gls{LIDAR} 3D. Las cámaras entregan mayor información visual y permiten
incorporar reconocimiento de objetos, pero aumentan la carga de procesamiento y
pueden verse afectadas por las condiciones de iluminación~\citep{tourani2022visual}. Los sensores \gls{LIDAR} 3D ofrecen una representación más completa del entorno, porque no se limitan a un único plano de barrido, a cambio de un mayor volumen de datos que procesar. En robots de interiores, en cambio, el SLAM basado en \gls{LIDAR} 2D resulta más económico que su equivalente en 3D y constituye la alternativa predominante \citep{ran2025review2dlidar}.

En relación con el mapeo, existen diferentes soluciones de SLAM 2D,
entre ellas GMapping, Hector SLAM, Cartographer y
\texttt{slam\_toolbox}. Estas herramientas presentan diferencias en consumo
computacional, dependencia de la odometría, capacidad de cierre de lazo e
integración con \gls{ROS2}. La Tabla~\ref{tab:comparativa_slam} resume la
comparación general de estas alternativas.

\begin{table}[htbp]
\centering
\caption{Comparación general de alternativas para SLAM 2D.}
\label{tab:comparativa_slam}
\small
\begin{tabular}{P{2.4cm} P{5.3cm} P{5.9cm}}
\toprule
\textbf{Alternativa} &
\textbf{Fortaleza principal} &
\textbf{Limitación principal} \\
\midrule

GMapping & Técnica de mapeo por rejilla con filtro de partículas
\citep{grisetti2007gmapping}, con buen desempeño en tiempo real y consumo de
memoria acotado en escenas pequeñas \citep{ran2025review2dlidar}.
& Depende de una odometría consistente y carece de detección de cierre de lazo,
y su portabilidad a \gls{ROS2} es menor que la de las alternativas más recientes
\citep{ran2025review2dlidar,ngoc2023efficient,macenski2021slamtoolbox}. \\

Hector SLAM & Puede operar con baja dependencia de la odometría de ruedas
\citep{kohlbrecher2011hector}. & Requiere un \gls{LIDAR} de alta tasa de barrido
para contener la deriva, y presenta deriva en el cierre de lazo en escenas
extensas \citep{ran2025review2dlidar}. \\

Cartographer &
Buena consistencia global y capacidad de cierre de lazo
\citep{hess2016cartographer}. &
Mayor complejidad de configuración y ajuste
\citep{ngoc2023efficient,macenski2023desks}. En una comparación independiente de
cinco algoritmos SLAM 2D, Cartographer presentó el mayor uso medio de CPU
\citep{trejos2022slam}. \\

\texttt{slam\_toolbox} &
Integración con \gls{ROS2}, grafos de poses y almacenamiento de mapas. &
En modo asíncrono puede omitir mediciones si el procesamiento de la anterior se retrasa
\citep{macenski2021slamtoolbox}. \\

\bottomrule
\end{tabular}
\end{table}

En consecuencia, la elección del sensor de percepción y de la herramienta de
SLAM depende del costo, la carga computacional, la dependencia de la odometría
y el grado de integración con \gls{ROS2}.

\section{Navegación y exploración autónoma}
\label{sec:exploracion}

La navegación autónoma combina dos decisiones de diseño: qué planificador
global calcula la ruta y qué controlador local la sigue evitando obstáculos.
Las alternativas disponibles para cada una difieren en la representación
del entorno que requieren, en su costo computacional y en su respuesta ante
obstáculos dinámicos.

En la planificación global, los métodos de búsqueda sobre rejilla, como
Dijkstra y A$^*$ \citep[Sec.~2.2.2]{lavalle2006planning}, siguen siendo los
más utilizados en mapas de ocupación 2D, porque operan directamente sobre la
rejilla y permiten replanificar de forma periódica con bajo costo
computacional \citep{macenski2023desks}. Los métodos basados en muestreo,
como RRT y RRT$^*$, se emplean sobre todo en espacios continuos de mayor
dimensión o en sistemas con restricciones de movimiento
\citep{lavalle2006planning,karaman2011optimal}. Los métodos de optimización de trayectorias
mejoran la suavidad y la factibilidad de la ruta, a cambio de un mayor costo
de cómputo \citep{macenski2023desks}.

En el control local coexisten tres familias. Los métodos reactivos, como el
enfoque de ventana dinámica, evalúan un conjunto reducido de velocidades
alcanzables y son habituales en robots diferenciales por su bajo costo y su
buena respuesta ante obstáculos cercanos. Los métodos geométricos, como Pure
Pursuit, siguen la ruta con menor esfuerzo de cálculo, pero responden peor
ante obstáculos dinámicos. Los métodos predictivos optimizan una secuencia
futura de comandos considerando la dinámica del robot, a costa de una mayor
carga computacional \citep{macenski2023desks}.

\subsection{Plataformas para navegación autónoma}
\label{sec:plataformas_nav}

Los métodos de planificación y control descritos no se ejecutan de forma
aislada, sino a través de una plataforma que los integra y coordina. Para esto
existen distintas alternativas, que van desde frameworks de navegación ya
consolidados hasta arquitecturas desarrolladas de forma específica.
Por su parte, Nav2 es el framework de navegación de referencia en \gls{ROS2}. Su estructura
modular permite combinar planificadores como NavFn y Smac con controladores
como DWB, Regulated Pure Pursuit y MPPI, y orquestar la navegación mediante
árboles de comportamiento \citep{nav2_doc}. Su predecesor en ROS~1,
\texttt{move\_base}, cumplió un rol equivalente durante años, pero quedó ligado
a una distribución cuyo soporte ya finalizó \citep{ros_noetic_eol}. En un
ámbito distinto, Autoware surgió como una pila de navegación orientada a
vehículos autónomos \citep{kato2018autoware}. La documentación oficial actual
lo describe como una pila integral de conducción autónoma que cubre desde el
sensado hasta el control \citep{autoware2026docs}. Por ello, su ámbito de
aplicación y sus requerimientos de percepción y cómputo siguen siendo
distintos a los de una plataforma de interiores como la desarrollada en este
trabajo. También existen soluciones propietarias integradas en robots
comerciales, cuyo funcionamiento interno no siempre es accesible.

Como alternativa a estos frameworks, es posible desarrollar una arquitectura
propia, apoyándose en bibliotecas de planificación de movimiento como OMPL
\citep{sucan2012ompl,ompl2026docs} o integrando los componentes mediante otros
middleware robóticos. Los algoritmos
de planificación y control pueden ser equivalentes, pero cambian las
interfaces, las herramientas de supervisión y los mecanismos de comunicación.
En términos generales, un framework consolidado reduce el esfuerzo de
integración y aporta componentes ya probados, mientras que una arquitectura
propia ofrece mayor control a cambio de un mayor esfuerzo de desarrollo.

\subsection{Exploración autónoma}
\label{sec:exploracion_alternativas}

Las estrategias de exploración pueden agruparse en enfoques basados en
fronteras, métodos guiados por ganancia de información y propuestas de
\textit{active SLAM}.
La exploración basada en fronteras, propuesta originalmente por
\citet{yamauchi1997frontier}, utiliza como candidatos las regiones que
separan el espacio conocido de las zonas todavía no observadas. Herramientas
como \texttt{explore\_lite} representan ejemplos de este enfoque, incluido
su port a \gls{ROS2} \citep{horner_mexplore,alvarez2021mexplore_ros2}. Su
principal ventaja es la facilidad de integración con mapas de ocupación y
sistemas de navegación ya existentes.

Los métodos basados en ganancia de información seleccionan metas estimando
cuánto conocimiento adicional podría obtenerse desde cada posición. Estos
métodos pueden considerar visibilidad, superficie esperada o reducción de
incertidumbre, pero requieren modelos y cálculos adicionales.

Las estrategias de \textit{active SLAM} incorporan simultáneamente la
cobertura del entorno, la incertidumbre de localización y la consistencia del
mapa. En desarrollos recientes también se investigan métodos basados en
muestreo, aprendizaje y coordinación de múltiples robots
\citep{placed2023active}.
Los enfoques basados en fronteras continúan siendo utilizados por su menor
complejidad y facilidad de integración. En cambio, los métodos de ganancia de
información y \textit{active SLAM} pueden generar decisiones más informadas,
aunque demandan una representación más completa del estado del sistema y una
mayor capacidad de procesamiento.

\section{Simulación de robots móviles}
\label{sec:gemelo_digital}

Entre las alternativas disponibles se encuentran simuladores de propósito
general, como Gazebo Classic, Webots y CoppeliaSim, además de plataformas orientadas
a una mayor fidelidad visual y física, como Isaac Sim. Estas herramientas
difieren en la precisión de sus modelos físicos, la representación de
sensores, la integración con plataformas robóticas, los recursos
computacionales requeridos y la facilidad para construir escenarios de
prueba.
Gazebo Classic permite simular robots y entornos tridimensionales mediante
motores de física, sensores y extensiones basadas en plugins
\citep{osrf2025gazeboclassic}. El Gazebo actual corresponde a la línea
sucesora mantenida por el proyecto y se distribuye separadamente de Gazebo
Classic \citep{openrobotics2024gazebo}. Webots y CoppeliaSim, presentados
originalmente en \citet{michel2004webots} y \citet{rohmer2013coppeliasim},
continúan disponibles como simuladores robóticos con modelado de robots y
entornos tridimensionales, de acuerdo con su documentación oficial reciente
\citep{cyberbotics2025webots,coppelia2025version}. Isaac Sim incorpora
capacidades avanzadas de representación gráfica, generación de
datos sintéticos y simulación de sensores, aunque requiere una mayor
capacidad de procesamiento \citep{nvidia2025isaacsim}.

Conviene precisar la nomenclatura, porque el nombre Gazebo designa
actualmente a dos programas distintos. La familia de versiones numeradas,
cuya última entrega es la 11, pasó a denominarse Gazebo Classic y alcanzó
su fin de soporte en enero de 2025 \citep{osrf2025gazeboclassic}. Su
sucesor, desarrollado inicialmente bajo el nombre Ignition y con versiones
identificadas por nombre en lugar de número, es el que hoy se distribuye
simplemente como Gazebo \citep{openrobotics2025gzmigration}.

La selección de un simulador depende del propósito de la prueba. Para validar
la cinemática, los marcos de referencia y la comunicación entre componentes,
puede ser suficiente un entorno de complejidad moderada. En cambio, las
aplicaciones que requieren percepción visual, interacción física detallada o
entrenamiento de sistemas de aprendizaje pueden necesitar simuladores con
mayor fidelidad.

\section*{Conclusión del capítulo}

La revisión muestra que las soluciones actuales para robótica móvil se
orientan hacia arquitecturas modulares basadas en \gls{ROS2}, percepción y
mapeo del entorno, planificación y exploración autónomas, y simuladores de
mayor fidelidad. Estas alternativas amplían las capacidades del sistema, pero suelen aumentar el costo, el consumo de recursos y la complejidad de
implementación. La selección concreta de tecnologías para el prototipo, que
busca un equilibrio entre estas capacidades y las restricciones de una
plataforma de bajo costo, se presenta en el Capítulo~\ref{ch:solucion}.